\documentclass[10pt,a4paper]{amsart}
\usepackage[margin=19mm]{geometry}
\usepackage{amsmath,amssymb,mathtools}
\usepackage[scr=rsfs,cal=euler]{mathalfa}
\usepackage{xfrac}
\usepackage[unicode,hidelinks,bookmarks=false]{hyperref}
\newcommand{\ie}{i.e.\ }
\newcommand{\cf}{cf.\ }
\newcommand{\resp}{respectively\ }
\newcommand{\Subsection}{\subsection}
\providecommand{\nobreakdash}{\nobreak\hbox{-}}
\setlength{\emergencystretch}{2em}
\multlinegap0pt
\usepackage{amssymb}
\newtheorem{lem}{Lemma}
\def\R{\mathbb{R}}
\def\T{\mathcal{T}}
\title{From Triangular Array Progression\\ to Near-Log-Concave State Transitions}
\author{Levent Ali Meng\"ut\"urk}
\date{Source: arXiv:2608.03956v5 (CC BY 4.0)}
\begin{document}
\maketitle

\noindent\emph{Source and licence.} From Triangular Array Progression\\ to Near-Log-Concave State Transitions. Version arXiv:2608.03956v5 (CC BY 4.0), distributed by arXiv under the Creative Commons Attribution 4.0 International licence, http://creativecommons.org/licenses/by/4.0/, as recorded in the arXiv licence metadata for this exact version and shown on the abstract page https://arxiv.org/abs/2608.03956v5. Full licence notice: \texttt{licenses/arxiv-2608.03956-CC-BY-4.0.txt}.\par\smallskip

\noindent\textit{Editorial scope.} Counted unit: the article's lem mainlemmaone (source0.tex lines 180--187). The article's complete original proof of it is delivered with it (source0.tex lines 188--245, one \texttt{proof} environment of 58 lines). No mathematics is written, repaired or re-derived: the statement and the proof are the article's own lines, in the article's own notation, and no retained line is substituted or re-flowed.  Retained uncounted as the source-local context the proof needs: lines 149--151.  Nothing was omitted from any retained span, so the package carries no omission record.  No retained line contains a citation, so the wrapper carries no bibliography and no bibliography entry.  No retained line contains a figure environment, an image-inclusion command or drawing code, so no image is delivered and the package declares no asset dependency.  Editorial formatting only, in addition: an AMS \texttt{amsart} preamble that restates the article's own declarations and macros used by the retained lines, namely the environments \texttt{lem} on separate counters, together with the standard packages those lines need.  The retained environments print in this document as Lemma 1 in order of appearance, whereas the article prints the same items with the numbering of its own full document.  The complete native source of the article as distributed in the arXiv e-print for this exact version is bundled unchanged as \texttt{sources/206799/source0.tex}.
\begin{align}
&\psi(j) = \sum_{k=1}^{j-1} \psi(k) \hspace{0.1in} \text{for $j \geq 4$}. \label{columnprogressionfunc}
\end{align}

\begin{lem}
\label{mainlemmaone}
Define an infinite sub-triangular array $\tau$ by
\begin{align}
\tau_{n,j} \triangleq \T_{n + 3, j + 3} \notag
\end{align}
for every $n \geq 0$ and $1 \leq j \leq n +1$. Then, $\tau_{n,*}$ is a strictly-positive log-concave sequence for every $n\geq 2$.
\end{lem}

\begin{proof}
To begin with, every element in $\T$ is strictly-positive, hence, every element in $\tau$ is strictly-positive. Compute the first four sequences of $\tau$, which yield 
\begin{align}
\tau_{0,*} = [1] \hspace{0.1in} \text{and} \hspace{0.1in} \tau_{1,*} = [5,1] \hspace{0.1in} \text{and} \hspace{0.1in} \tau_{2,*} = [8,10,1] \hspace{0.1in} \text{and} \hspace{0.1in} \tau_{3,*} = [11,16,20,1], \label{initialtauseq} 
\end{align}
for which log-concavity is not relevant for $\tau_{0,*}$ nor $\tau_{1,*}$, whereby $\tau_{2,*}$ and $\tau_{3,*}$ are log-concave sequences, since they satisfy 
\begin{align}
\tau_{2,1}\tau_{2,3} \leq \tau_{2,2}^2 \hspace{0.1in} \text{and} \hspace{0.1in} \tau_{3,1}\tau_{3,3} \leq \tau_{3,2}^2 \hspace{0.1in} \text{and} \hspace{0.1in} \tau_{3,2}\tau_{3,4} \leq \tau_{3,3}^2, \label{initiallogconcaveprop} 
\end{align}
respectively. Define the following continuous functions: 
\begin{align}
f_1(x) = 11 + 3x, \,\,\,\, f_2(x) = 16 + 6x, \,\,\,\, f_3(x) = 20 + 12x, \notag 
\end{align}
for any $x\in\R_+$. Since $\T_{n,j} = \T_{n-1,j} + \psi(j)$ for every $n \geq 1$ and $1 \leq j \leq n-1$, the first three columns of $\tau$ define discrete points on $f_1$, $f_2$ and $f_3$, respectively. Note also that we have the positivity:
\begin{align}
f_2(x)^2 - f_1(x)f_3(x) = 36, \notag 
\end{align}
for any $x\in\R_+$. This implies that any sequence formed from the first three columns of $\tau$ satisfy log-concavity:
\begin{align}
\label{funcapprox}
\tau_{n,1}\tau_{n,3} \leq \tau_{n,2}^2,
\end{align}
for every $n\geq 2$. Next, from (\ref{columnprogressionfunc}), we can write
\begin{equation}
\label{columnprogressiontwo}
\psi(j) =
\begin{cases}
j -1 & \text{for $1 \leq j \leq 4$},\\
2\psi(j-1) &  \text{for $j \geq 5$}. \\
\end{cases}
\end{equation}
In addition, since $\T_{2,2}$ = 2, $\T_{3,3} = 3$ and $\T_{n,n} = \sum_{k=1}^{n-2} \T_{n-k,n-k}$ for every $n \geq 4$, we also have the following diagonal relation:
\begin{align}
\T_{n,n} = 2\T_{n-1,n-1} \label{remakeofdiagonals}
\end{align}
for every $ n \geq 5$. Therefore, using (\ref{initialtauseq}), (\ref{columnprogressiontwo}), (\ref{remakeofdiagonals}) and $\T_{n,j} = \T_{n-1,j} + \psi(j)$ for every $n \geq 1$ and $1 \leq j \leq n-1$, we must have
\begin{align}
\tau_{n,j} = 2 \tau_{n-1,j-1}, \label{maindiagonalrelationfirsteq}
\end{align}
for every $n \geq 2$ and $2 \leq j \leq n$. Moreover, since we have $\T_{n,n+1} = 1$ for every $n \geq 0$, we also have $\tau_{n,n+1} = 1$, which further means $\tau_{n,n+1} = \tau_{n-1,n}$ for every $n \geq 1$. Accordingly, we can write
\begin{align}
\label{maindiagonalrelation}
\tau_{n,j} = \alpha_j(n) \tau_{n-1,j-1}
\end{align}
for every $n \geq 2$ and $2 \leq j \leq n+1$, where we defined the coefficient
\begin{equation}
\alpha_j(n) =
\begin{cases}
1 & \text{for $j = n+1$},\\
2 &  \text{for $2 \leq j \leq n$}. \notag
\end{cases}
\end{equation}
Thus, using (\ref{funcapprox}) and (\ref{maindiagonalrelation}), we must have 
\begin{align}
\tau_{n,j-1}\tau_{n,j+1} \leq \tau_{n,j}^2, \notag 
\end{align}
for every $n\geq 2$ and $2 \leq j \leq n$. Therefore, the sequence $\tau_{n,*}$ is a strictly-positive log-concave sequence for every $n\geq 2$.
\end{proof}

\end{document}
